% ----- auto-generated from week_06/Week*.html via pandoc -----
\setchapterimage[6cm]{}
\setchapterpreamble[u]{\margintoc}
\chapter{Backstepping Control}
\labch{week_06}

\section{Week 6: Backstepping Control}\label{week-6-backstepping-control}}

Recursive Lyapunov Design for Nonlinear Systems

M.Sc. Control for Robots --- University of West London --- AY 2025--26\\
Duration: 3 Hours (Lecture + Hands-on Lab)

\subsection{Section 0 --- Week 5 Recall: Robust Control --- Exam Practice (20 min)}\label{section-0-week-5-recall-robust-control-exam-practice-20-min}}

\subsubsection{Question 1 --- Robust Sliding-Mode Analysis}\label{question-1-robust-sliding-mode-analysis}}

Consider a manipulator with model uncertainty:

\textbackslash{[}
M(q)\textbackslash ddot\{q\} + C(q,\textbackslash dot\{q\})\textbackslash dot\{q\} + G(q) = \textbackslash tau + d(t), \textbackslash qquad \textbackslash\textbar d(t)\textbackslash\textbar{} \textbackslash le d\_\{\textbackslash max\}
\textbackslash{]}

The robust controller is:

\textbackslash{[}
\textbackslash tau = \textbackslash hat\{M\}(q)\textbackslash bigl{[}\textbackslash ddot\{q\}\_d + K\_D \textbackslash dot\{e\} + K\_P e\textbackslash bigr{]} + \textbackslash hat\{C\}\textbackslash dot\{q\} + \textbackslash hat\{G\} + \textbackslash tau\_\{\textbackslash text\{robust\}\}
\textbackslash{]}

with the robust term

\textbackslash{[}
\textbackslash tau\_\{\textbackslash text\{robust\}\} = \textbackslash rho(t)\textbackslash,\textbackslash frac\{s\}\{\textbackslash\textbar s\textbackslash\textbar\}, \textbackslash qquad s = \textbackslash dot\{e\} + \textbackslash Lambda\textbackslash, e
\textbackslash{]}

where \textbackslash(e = q\_d - q\textbackslash), \textbackslash(\textbackslash dot\{e\} = \textbackslash dot\{q\}\_d - \textbackslash dot\{q\}\textbackslash), and \textbackslash(\textbackslash Lambda \textbackslash succ 0\textbackslash).

\paragraph{Model Answer --- Closed-Loop Error Dynamics}\label{model-answer-closed-loop-error-dynamics}}

\textbf{Step 1 --- Substitution.} Write the true dynamics as
\textbackslash(M\textbackslash ddot\{q\} = \textbackslash tau + d - C\textbackslash dot\{q\} - G\textbackslash).
Substitute the controller:

\textbackslash{[}
M\textbackslash ddot\{q\} = \textbackslash hat\{M\}(\textbackslash ddot\{q\}\_d + K\_D\textbackslash dot\{e\} + K\_P e) + \textbackslash hat\{C\}\textbackslash dot\{q\} + \textbackslash hat\{G\}
+ \textbackslash rho\textbackslash frac\{s\}\{\textbackslash\textbar s\textbackslash\textbar\} + d - C\textbackslash dot\{q\} - G
\textbackslash{]}

Define the model errors
\textbackslash(\textbackslash tilde\{M\} = M - \textbackslash hat\{M\}\textbackslash), \textbackslash(\textbackslash tilde\{C\} = C - \textbackslash hat\{C\}\textbackslash), \textbackslash(\textbackslash tilde\{G\} = G - \textbackslash hat\{G\}\textbackslash),
and let the lumped uncertainty be

\textbackslash{[}
\textbackslash eta(t) = \textbackslash tilde\{M\}\textbackslash ddot\{q\} + \textbackslash tilde\{C\}\textbackslash dot\{q\} + \textbackslash tilde\{G\} - d(t)
\textbackslash{]}

\textbf{Step 2 --- Error dynamics.}
With \textbackslash(\textbackslash ddot\{e\} = \textbackslash ddot\{q\}\_d - \textbackslash ddot\{q\}\textbackslash) and after algebraic manipulation:

\textbackslash{[}
M\textbackslash ddot\{e\} + (C + M K\_D)\textbackslash dot\{e\} + M K\_P e = -\textbackslash eta(t) + \textbackslash rho\textbackslash frac\{s\}\{\textbackslash\textbar s\textbackslash\textbar\}
\textbackslash{]}

In terms of the composite variable \textbackslash(s = \textbackslash dot\{e\} + \textbackslash Lambda e\textbackslash):

\textbackslash{[}
M\textbackslash dot\{s\} = -Cs + \textbackslash eta(t) + \textbackslash rho\textbackslash frac\{s\}\{\textbackslash\textbar s\textbackslash\textbar\}
\textbackslash{]}

(using the choice \textbackslash(K\_D = \textbackslash Lambda\textbackslash), \textbackslash(K\_P = \textbackslash dot\{\textbackslash Lambda\}\textbackslash) to simplify).

\textbf{Step 3 --- Lyapunov analysis.} Choose
\textbackslash(V = \textbackslash tfrac\{1\}\{2\}\textbackslash,s\^{}T M\textbackslash, s\textbackslash). Then:

\textbackslash{[}
\textbackslash dot\{V\} = s\^{}T M\textbackslash dot\{s\} + \textbackslash tfrac\{1\}\{2\}\textbackslash,s\^{}T\textbackslash dot\{M\}\textbackslash,s
= s\^{}T\textbackslash!\textbackslash bigl{[}-Cs + \textbackslash eta + \textbackslash rho\textbackslash tfrac\{s\}\{\textbackslash\textbar s\textbackslash\textbar\}\textbackslash bigr{]} + \textbackslash tfrac\{1\}\{2\}\textbackslash,s\^{}T\textbackslash dot\{M\}\textbackslash,s
\textbackslash{]}

Using the skew-symmetry property \textbackslash(s\^{}T(\textbackslash dot\{M\} - 2C)s = 0\textbackslash):

\textbackslash{[}
\textbackslash dot\{V\} = s\^{}T\textbackslash eta + \textbackslash rho\textbackslash\textbar s\textbackslash\textbar{}
\textbackslash{]}

Since \textbackslash(\textbackslash tau\_\{\textbackslash text\{robust\}\}\textbackslash) opposes the uncertainty
(note the sign convention: \textbackslash(\textbackslash rho \textbackslash frac\{s\}\{\textbackslash\textbar s\textbackslash\textbar\}\textbackslash) acts in the direction of \textbackslash(s\textbackslash), so the correct robust term to \emph{counteract} uncertainty is \textbackslash(\textbackslash tau\_\{\textbackslash text\{robust\}\} = -\textbackslash rho \textbackslash frac\{s\}\{\textbackslash\textbar s\textbackslash\textbar\}\textbackslash)). Re-doing with the standard sign:

\textbackslash{[}
\textbackslash dot\{V\} = s\^{}T\textbackslash eta - \textbackslash rho\textbackslash\textbar s\textbackslash\textbar{} \textbackslash le \textbackslash\textbar\textbackslash eta\textbackslash\textbar\textbackslash,\textbackslash\textbar s\textbackslash\textbar{} - \textbackslash rho\textbackslash\textbar s\textbackslash\textbar{} = (\textbackslash\textbar\textbackslash eta\textbackslash\textbar{} - \textbackslash rho)\textbackslash\textbar s\textbackslash\textbar{}
\textbackslash{]}

\textbf{Conclusion:} If \textbackslash(\textbackslash rho(t) \textbackslash ge \textbackslash\textbar\textbackslash eta(t)\textbackslash\textbar\textbackslash), then \textbackslash(\textbackslash dot\{V\} \textbackslash le 0\textbackslash) and the tracking error \textbackslash(s \textbackslash to 0\textbackslash). The sliding variable converges, guaranteeing bounded tracking despite uncertainty. By Barbalat\textquotesingle s lemma, \textbackslash(s(t) \textbackslash to 0\textbackslash) asymptotically, which implies \textbackslash(e(t) \textbackslash to 0\textbackslash) exponentially since \textbackslash(\textbackslash dot\{e\} = s - \textbackslash Lambda e\textbackslash) is a stable filter.

\subsubsection{Question 2 --- Robust vs. Adaptive: Mathematical Comparison}\label{question-2-robust-vs.-adaptive-mathematical-comparison}}

\paragraph{Model Answer}\label{model-answer}}

\textbf{Adaptive control} uses an augmented Lyapunov function:

\textbackslash{[}
V\_\{\textbackslash text\{adaptive\}\} = \textbackslash frac\{1\}\{2\}\textbackslash,s\^{}T M\textbackslash, s + \textbackslash frac\{1\}\{2\}\textbackslash,\textbackslash tilde\{\textbackslash theta\}\^{}T \textbackslash Gamma\^{}\{-1\} \textbackslash tilde\{\textbackslash theta\}
\textbackslash{]}

where \textbackslash(\textbackslash tilde\{\textbackslash theta\} = \textbackslash theta - \textbackslash hat\{\textbackslash theta\}\textbackslash) is the parameter estimation error
and \textbackslash(\textbackslash Gamma \textbackslash succ 0\textbackslash) is the adaptation gain matrix. The adaptation law
\textbackslash(\textbackslash dot\{\textbackslash hat\{\textbackslash theta\}\} = \textbackslash Gamma\textbackslash, Y\^{}T s\textbackslash) is designed so that
the cross-term \textbackslash(s\^{}T Y \textbackslash tilde\{\textbackslash theta\}\textbackslash) in \textbackslash(\textbackslash dot\{V\}\textbackslash) cancels exactly with
\textbackslash(\textbackslash tilde\{\textbackslash theta\}\^{}T \textbackslash Gamma\^{}\{-1\}\textbackslash dot\{\textbackslash hat\{\textbackslash theta\}\}\textbackslash), yielding:

\textbackslash{[}
\textbackslash dot\{V\}\_\{\textbackslash text\{adaptive\}\} = -s\^{}T K\_D\textbackslash, s \textbackslash le 0
\textbackslash{]}

\textbf{Robust control} uses a simpler Lyapunov function:

\textbackslash{[}
V\_\{\textbackslash text\{robust\}\} = \textbackslash frac\{1\}\{2\}\textbackslash,s\^{}T M\textbackslash, s
\textbackslash{]}

No parameter estimation is required because the robust term directly dominates
the uncertainty via the bound \textbackslash(\textbackslash rho \textbackslash ge \textbackslash\textbar\textbackslash eta\textbackslash\textbar\textbackslash):

\textbackslash{[}
\textbackslash dot\{V\}\_\{\textbackslash text\{robust\}\} \textbackslash le -(\textbackslash rho - \textbackslash\textbar\textbackslash eta\textbackslash\textbar)\textbackslash\textbar s\textbackslash\textbar{} \textbackslash le 0
\textbackslash{]}

\textbf{Why robust control does not need parameter estimation:}
Instead of identifying \emph{what} the uncertainty is (as adaptive control does), robust control only needs to know \emph{how large} the uncertainty can be. It overpowers the worst case with a dominating gain \textbackslash(\textbackslash rho\textbackslash).

\textbf{Trade-off:}

\begin{longtable}[]{@{}lll@{}}
\toprule\noalign{}
Aspect & Adaptive Control & Robust Control \\
\midrule\noalign{}
\endhead
\bottomrule\noalign{}
\endlastfoot
Uncertainty handling & Estimates and cancels parametric uncertainty & Overpowers bounded uncertainty without identification \\
Conservatism & Low --- adapts to actual parameters & High --- must use worst-case bound \\
Control effort & Decreases as estimates converge & Remains high (constant worst-case gain) \\
Chattering & None (smooth update law) & Possible with discontinuous sign function \\
Unstructured uncertainty & Cannot handle (requires linear parameterisation) & Handles any bounded disturbance \\
Convergence & Asymptotic (with PE: exponential) & Finite-time to sliding surface \\
\end{longtable}

\subsection{Section 1 --- Learning Outcomes}\label{section-1-learning-outcomes}}

\begin{itemize}
\tightlist
\item
  Understand the backstepping design methodology for strict-feedback nonlinear systems
\item
  Derive a backstepping controller step-by-step using recursive Lyapunov design
\item
  Apply backstepping to manipulator, mobile robot, and drone systems
\item
  Implement and simulate backstepping controllers in MATLAB
\item
  Compare backstepping with computed torque, MPC, adaptive, and robust control
\end{itemize}

\subsection{Section 2 --- Session Roadmap (3 Hours)}\label{section-2-session-roadmap-3-hours}}

0:00 -- 0:20

\paragraph{Week 5 Recall --- Robust Control Exam Practice}\label{week-5-recall-robust-control-exam-practice}}

Solve two exam-style questions on robust control, Lyapunov analysis, and comparison with adaptive control.

0:20 -- 1:00

\paragraph{Part 1 --- Backstepping Theory}\label{part-1-backstepping-theory}}

Strict-feedback systems, virtual control, recursive Lyapunov design, n-step generalisation.

1:00 -- 1:30

\paragraph{Part 2 --- Backstepping for Manipulator}\label{part-2-backstepping-for-manipulator}}

Convert Euler-Lagrange dynamics to strict-feedback form, design backstepping torque controller, full Lyapunov proof, MATLAB code.

1:30 -- 2:00

\paragraph{Part 3 --- Backstepping for Mobile Robot}\label{part-3-backstepping-for-mobile-robot}}

Chained-form unicycle model, backstepping position tracking design, MATLAB code.

2:00 -- 2:30

\paragraph{Part 4 --- Backstepping for Drone}\label{part-4-backstepping-for-drone}}

Cascade position--attitude structure, multi-loop backstepping, MATLAB code.

2:30 -- 3:00

\paragraph{Hands-on Activities, Quiz \& Summary}\label{hands-on-activities-quiz-summary}}

Gain tuning experiments, comparative study, quiz, and method comparison table.

\subsection{Section 3 --- Part 1: Backstepping Theory}\label{section-3-part-1-backstepping-theory}}

\subsubsection{3.1 Motivation --- Why Backstepping?}\label{motivation-why-backstepping}}

In previous weeks, we studied controllers that either require exact cancellation of nonlinearities
(computed torque), predict future behaviour (MPC), estimate unknown parameters on-line (adaptive),
or overpower bounded uncertainties (robust). \textbf{Backstepping} takes a fundamentally different
approach: it exploits the \emph{triangular (cascade) structure} of the system to build a Lyapunov
function and controller \emph{one subsystem at a time}, recursively.

\textbf{Key Insight:} Many physical systems --- from robotic manipulators to quadrotors ---
have a \emph{strict-feedback} (cascade) structure where one state acts as the "input" to the previous
subsystem. Backstepping systematically exploits this structure.

\textbf{Figure 1:} Backstepping cascade structure. The design proceeds from the outermost subsystem (Step 1) back to the actual control input (Step n).

\subsubsection{3.2 Strict-Feedback Form}\label{strict-feedback-form}}

A nonlinear system is in \textbf{strict-feedback form} if it can be written as:

\textbackslash{[}
\textbackslash begin\{aligned\}
\textbackslash dot\{x\}\_1 \&= f\_1(x\_1) + g\_1(x\_1)\textbackslash,x\_2 \textbackslash\textbackslash{}
\textbackslash dot\{x\}\_2 \&= f\_2(x\_1, x\_2) + g\_2(x\_1, x\_2)\textbackslash,x\_3 \textbackslash\textbackslash{}
\&\textbackslash;\textbackslash;\textbackslash vdots \textbackslash\textbackslash{}
\textbackslash dot\{x\}\_\{n-1\} \&= f\_\{n-1\}(x\_1,\textbackslash ldots,x\_\{n-1\}) + g\_\{n-1\}(x\_1,\textbackslash ldots,x\_\{n-1\})\textbackslash,x\_n \textbackslash\textbackslash{}
\textbackslash dot\{x\}\_n \&= f\_n(x\_1,\textbackslash ldots,x\_n) + g\_n(x\_1,\textbackslash ldots,x\_n)\textbackslash,u
\textbackslash end\{aligned\}
\textbackslash{]}

Key properties:

\begin{itemize}
\tightlist
\item
  Each equation depends on all previous states plus \emph{one} new state \textbackslash(x\_\{i+1\}\textbackslash).
\item
  The new state \textbackslash(x\_\{i+1\}\textbackslash) enters \emph{multiplicatively} through \textbackslash(g\_i(\textbackslash cdot)\textbackslash).
\item
  Only the \emph{last} equation contains the true control input \textbackslash(u\textbackslash).
\item
  \textbackslash(g\_i(\textbackslash cdot) \textbackslash neq 0\textbackslash) everywhere (controllability condition).
\end{itemize}

\textbf{Figure 2:} Triangular dependency structure of strict-feedback form. Each row depends only on states from current and previous subsystems.

\subsubsection{3.3 Two-Step Backstepping --- Complete Derivation}\label{two-step-backstepping-complete-derivation}}

Consider the simplest strict-feedback system (two states, one input):

\textbackslash{[}
\textbackslash begin\{aligned\}
\textbackslash dot\{x\}\_1 \&= f\_1(x\_1) + g\_1(x\_1)\textbackslash,x\_2 \textbackslash\textbackslash{}
\textbackslash dot\{x\}\_2 \&= f\_2(x\_1,x\_2) + g\_2(x\_1,x\_2)\textbackslash,u
\textbackslash end\{aligned\}
\textbackslash{]}

\textbf{Goal:} Drive \textbackslash(x\_1 \textbackslash to x\_\{1,\textbackslash text\{ref\}\}\textbackslash) (a desired reference).

\paragraph{Step 1 --- Stabilise the \textbackslash(x\_1\textbackslash) Subsystem (Virtual Control Design)}\label{step-1-stabilise-the-x_1-subsystem-virtual-control-design}}

Define the tracking error:

\textbackslash{[}
z\_1 = x\_1 - x\_\{1,\textbackslash text\{ref\}\}
\textbackslash{]}

Its dynamics are:

\textbackslash{[}
\textbackslash dot\{z\}\_1 = \textbackslash dot\{x\}\_1 - \textbackslash dot\{x\}\_\{1,\textbackslash text\{ref\}\} = f\_1(x\_1) + g\_1(x\_1)\textbackslash,x\_2 - \textbackslash dot\{x\}\_\{1,\textbackslash text\{ref\}\}
\textbackslash{]}

Imagine that \textbackslash(x\_2\textbackslash) is a control input that we can freely choose. We call this the
\textbf{virtual control} and design it as \textbackslash(x\_2 = \textbackslash alpha\_1(x\_1)\textbackslash) such that \textbackslash(z\_1 \textbackslash to 0\textbackslash).
Choose the Lyapunov function candidate:

\textbackslash{[}
V\_1 = \textbackslash frac\{1\}\{2\}\textbackslash,z\_1\^{}2
\textbackslash{]}

Its time derivative is:

\textbackslash{[}
\textbackslash dot\{V\}\_1 = z\_1\textbackslash,\textbackslash dot\{z\}\_1 = z\_1\textbackslash bigl{[}f\_1(x\_1) + g\_1(x\_1)\textbackslash,x\_2 - \textbackslash dot\{x\}\_\{1,\textbackslash text\{ref\}\}\textbackslash bigr{]}
\textbackslash{]}

If \textbackslash(x\_2\textbackslash) were our actual input, we would choose:

\textbackslash{[}
\textbackslash alpha\_1(x\_1) = \textbackslash frac\{1\}\{g\_1(x\_1)\}\textbackslash bigl{[}-f\_1(x\_1) + \textbackslash dot\{x\}\_\{1,\textbackslash text\{ref\}\} - k\_1\textbackslash,z\_1\textbackslash bigr{]}, \textbackslash qquad k\_1 \textgreater{} 0
\textbackslash{]}

This would give \textbackslash(\textbackslash dot\{V\}\_1 = -k\_1\textbackslash,z\_1\^{}2 \textbackslash le 0\textbackslash). But \textbackslash(x\_2\textbackslash) is not the actual input ---
it is a state. So we define the \textbf{backstepping error}:

\textbackslash{[}
z\_2 = x\_2 - \textbackslash alpha\_1(x\_1)
\textbackslash{]}

Substituting \textbackslash(x\_2 = z\_2 + \textbackslash alpha\_1\textbackslash) into \textbackslash(\textbackslash dot\{V\}\_1\textbackslash):

\textbackslash{[}
\textbackslash dot\{V\}\_1 = z\_1\textbackslash bigl{[}f\_1 + g\_1(z\_2 + \textbackslash alpha\_1) - \textbackslash dot\{x\}\_\{1,\textbackslash text\{ref\}\}\textbackslash bigr{]}
= -k\_1\textbackslash,z\_1\^{}2 + g\_1\textbackslash,z\_1\textbackslash,z\_2
\textbackslash{]}

The term \textbackslash(-k\_1 z\_1\^{}2\textbackslash) is stabilising; the coupling term \textbackslash(g\_1 z\_1 z\_2\textbackslash) must be handled in Step 2.

\paragraph{Step 2 --- Design the Actual Control \textbackslash(u\textbackslash) (Backstep)}\label{step-2-design-the-actual-control-u-backstep}}

Now we must drive \textbackslash(z\_2 \textbackslash to 0\textbackslash). Compute the dynamics of \textbackslash(z\_2\textbackslash):

\textbackslash{[}
\textbackslash dot\{z\}\_2 = \textbackslash dot\{x\}\_2 - \textbackslash dot\{\textbackslash alpha\}\_1
= f\_2(x\_1,x\_2) + g\_2(x\_1,x\_2)\textbackslash,u - \textbackslash frac\{\textbackslash partial \textbackslash alpha\_1\}\{\textbackslash partial x\_1\}\textbackslash dot\{x\}\_1
- \textbackslash frac\{\textbackslash partial \textbackslash alpha\_1\}\{\textbackslash partial x\_\{1,\textbackslash text\{ref\}\}\}\textbackslash dot\{x\}\_\{1,\textbackslash text\{ref\}\}
\textbackslash{]}

where we have used the chain rule on \textbackslash(\textbackslash alpha\_1(x\_1, x\_\{1,\textbackslash text\{ref\}\})\textbackslash). Augment the Lyapunov function:

\textbackslash{[}
V\_2 = V\_1 + \textbackslash frac\{1\}\{2\}\textbackslash,z\_2\^{}2 = \textbackslash frac\{1\}\{2\}\textbackslash,z\_1\^{}2 + \textbackslash frac\{1\}\{2\}\textbackslash,z\_2\^{}2
\textbackslash{]}

Its derivative is:

\textbackslash{[}
\textbackslash dot\{V\}\_2 = \textbackslash dot\{V\}\_1 + z\_2\textbackslash,\textbackslash dot\{z\}\_2
= -k\_1\textbackslash,z\_1\^{}2 + g\_1\textbackslash,z\_1\textbackslash,z\_2 + z\_2\textbackslash,\textbackslash dot\{z\}\_2
\textbackslash{]}

We want \textbackslash(\textbackslash dot\{V\}\_2 = -k\_1\textbackslash,z\_1\^{}2 - k\_2\textbackslash,z\_2\^{}2\textbackslash) for some \textbackslash(k\_2 \textgreater{} 0\textbackslash). This requires:

\textbackslash{[}
z\_2\textbackslash bigl{[}g\_1\textbackslash,z\_1 + \textbackslash dot\{z\}\_2\textbackslash bigr{]} = -k\_2\textbackslash,z\_2\^{}2
\textbackslash{]}
\textbackslash{[}
g\_1\textbackslash,z\_1 + f\_2 + g\_2\textbackslash,u - \textbackslash dot\{\textbackslash alpha\}\_1 = -k\_2\textbackslash,z\_2
\textbackslash{]}

Solving for the control input:

\textbackslash{[}
u = \textbackslash frac\{1\}\{g\_2(x\_1,x\_2)\}\textbackslash left{[}-f\_2(x\_1,x\_2) + \textbackslash dot\{\textbackslash alpha\}\_1
- g\_1(x\_1)\textbackslash,z\_1 - k\_2\textbackslash,z\_2\textbackslash right{]}
\textbackslash{]}

\textbf{Result:} With this control law,

\textbackslash{[}
\textbackslash dot\{V\}\_2 = -k\_1\textbackslash,z\_1\^{}2 - k\_2\textbackslash,z\_2\^{}2 \textbackslash le 0
\textbackslash{]}

which guarantees global asymptotic stability of \textbackslash((z\_1, z\_2) = (0,0)\textbackslash) by Lyapunov\textquotesingle s theorem.

\subsubsection{3.4 Generalisation to \textbackslash(n\textbackslash)-Step Backstepping}\label{generalisation-to-n-step-backstepping}}

For an \textbackslash(n\textbackslash)-dimensional strict-feedback system, the procedure generalises recursively:

Algorithm: n-Step Backstepping

For i = 1, 2, ..., n:
1. Define error: z\_i = x\_i - alpha\_\{i-1\} (where alpha\_0 = x\_\{1,ref\})
2. Compute dynamics: dz\_i/dt = f\_i + g\_i * x\_\{i+1\} - d(alpha\_\{i-1\})/dt
3. Augment Lyapunov: V\_i = V\_\{i-1\} + (1/2) * z\_i\^{}2
4. Compute dV\_i/dt and choose:
- If i \textless{} n: alpha\_i (virtual control) to make:
dV\_i/dt = -k\_1*z\_1\^{}2 - ... - k\_i*z\_i\^{}2 + g\_i*z\_i*z\_\{i+1\}
- If i = n: u (actual control) to make:
dV\_n/dt = -k\_1*z\_1\^{}2 - ... - k\_n*z\_n\^{}2
Result: Global asymptotic stability with V\_n as Lyapunov function.

\textbf{Figure 3:} Step-by-step backstepping design flowchart. Each step defines an error, augments the Lyapunov function, and designs a (virtual) control.

\subsubsection{3.5 Why Backstepping Works --- Intuition}\label{why-backstepping-works-intuition}}

\paragraph{Recursive Lyapunov Construction}\label{recursive-lyapunov-construction}}

At each step, we extend the Lyapunov function by adding \textbackslash(\textbackslash frac\{1\}\{2\}z\_i\^{}2\textbackslash).
Each new virtual control is designed to make the new terms in \textbackslash(\textbackslash dot\{V\}\textbackslash) negative definite
while leaving the previously designed terms unchanged.

\paragraph{"Backstepping" the Control}\label{backstepping-the-control}}

The name comes from the fact that we start from the outermost subsystem (where the reference
lives) and "step back" through the cascade, designing one layer at a time until we reach the
actual control input.

\paragraph{No Cancellation Required}\label{no-cancellation-required}}

Unlike feedback linearisation, backstepping does not cancel all nonlinearities. It only
cancels what is needed and uses the remaining structure constructively. This gives better
robustness and avoids wasting control effort.

\paragraph{Guaranteed Stability}\label{guaranteed-stability}}

The final Lyapunov function \textbackslash(V\_n = \textbackslash sum\_\{i=1\}\^{}n \textbackslash frac\{1\}\{2\}z\_i\^{}2\textbackslash) with
\textbackslash(\textbackslash dot\{V\}\_n = -\textbackslash sum\_\{i=1\}\^{}n k\_i z\_i\^{}2\textbackslash) provides a \emph{constructive} proof of
global asymptotic stability.

\textbf{Figure 5:} Lyapunov function augmentation. Each step adds a new error term, growing the "bowl" in state space while maintaining negative definiteness of V̇.

\paragraph{Think-Pair-Share (3 min)}\label{think-pair-share-3-min}}

\textbf{Question:} Why can\textquotesingle t we just use feedback linearisation for all nonlinear systems?
When does backstepping have an advantage?

\textbf{Discussion points:}

\begin{itemize}
\tightlist
\item
  Feedback linearisation requires an \emph{exact} model to cancel all nonlinearities ---
  any mismatch creates uncompensated dynamics.
\item
  Feedback linearisation may cancel \emph{useful} nonlinearities (e.g., natural damping).
\item
  Some systems are not feedback-linearisable (the relative degree conditions fail).
\item
  Backstepping provides a Lyapunov function at every step --- we can prove stability even
  with partial model knowledge.
\item
  Backstepping can be combined with adaptation (adaptive backstepping) to handle unknown
  parameters in the strict-feedback structure.
\end{itemize}

\subsection{Section 4 --- Part 2: Backstepping for Robot Manipulator}\label{section-4-part-2-backstepping-for-robot-manipulator}}

\subsubsection{4.1 Strict-Feedback Form of Manipulator Dynamics}\label{strict-feedback-form-of-manipulator-dynamics}}

The standard Euler-Lagrange dynamics of an \textbackslash(n\textbackslash)-DOF manipulator are:

\textbackslash{[}
M(q)\textbackslash ddot\{q\} + C(q,\textbackslash dot\{q\})\textbackslash dot\{q\} + G(q) = \textbackslash tau
\textbackslash{]}

Define the states \textbackslash(x\_1 = q\textbackslash) (joint positions) and \textbackslash(x\_2 = \textbackslash dot\{q\}\textbackslash) (joint velocities). Then:

\textbackslash{[}
\textbackslash begin\{aligned\}
\textbackslash dot\{x\}\_1 \&= x\_2 \textbackslash\textbackslash{}
\textbackslash dot\{x\}\_2 \&= M\^{}\{-1\}(x\_1)\textbackslash bigl{[}\textbackslash tau - C(x\_1,x\_2)\textbackslash,x\_2 - G(x\_1)\textbackslash bigr{]}
\textbackslash end\{aligned\}
\textbackslash{]}

This is a strict-feedback system with \textbackslash(f\_1 = 0\textbackslash), \textbackslash(g\_1 = I\textbackslash),
\textbackslash(f\_2 = -M\^{}\{-1\}(C x\_2 + G)\textbackslash), and \textbackslash(g\_2 = M\^{}\{-1\}\textbackslash).

\textbf{Figure 4:} 2-DOF planar manipulator and its decomposition into position and velocity subsystems for backstepping design.

\paragraph{Step 1 --- Virtual Control for Position Subsystem}\label{step-1-virtual-control-for-position-subsystem}}

Define the position error:

\textbackslash{[}
z\_1 = x\_1 - q\_d = q - q\_d
\textbackslash{]}

The dynamics of \textbackslash(z\_1\textbackslash) are:

\textbackslash{[}
\textbackslash dot\{z\}\_1 = \textbackslash dot\{q\} - \textbackslash dot\{q\}\_d = x\_2 - \textbackslash dot\{q\}\_d
\textbackslash{]}

Choose the Lyapunov function \textbackslash(V\_1 = \textbackslash frac\{1\}\{2\}\textbackslash,z\_1\^{}T z\_1\textbackslash). Its derivative:

\textbackslash{[}
\textbackslash dot\{V\}\_1 = z\_1\^{}T \textbackslash dot\{z\}\_1 = z\_1\^{}T(x\_2 - \textbackslash dot\{q\}\_d)
\textbackslash{]}

Design the virtual control \textbackslash(\textbackslash alpha\_1\textbackslash) (the desired velocity) so that if \textbackslash(x\_2 = \textbackslash alpha\_1\textbackslash), then \textbackslash(\textbackslash dot\{V\}\_1 \textless{} 0\textbackslash):

\textbackslash{[}
\textbackslash alpha\_1 = \textbackslash dot\{q\}\_d - K\_1\textbackslash,z\_1, \textbackslash qquad K\_1 = k\_1 I \textbackslash succ 0
\textbackslash{]}

This gives \textbackslash(\textbackslash dot\{V\}\_1 = -z\_1\^{}T K\_1 z\_1 + z\_1\^{}T z\_2\textbackslash) where \textbackslash(z\_2 = x\_2 - \textbackslash alpha\_1 = \textbackslash dot\{q\} - \textbackslash dot\{q\}\_d + K\_1 z\_1\textbackslash).

\paragraph{Step 2 --- Torque Control for Velocity Subsystem}\label{step-2-torque-control-for-velocity-subsystem}}

Define the velocity error:

\textbackslash{[}
z\_2 = x\_2 - \textbackslash alpha\_1 = \textbackslash dot\{q\} - \textbackslash dot\{q\}\_d + K\_1(q - q\_d)
\textbackslash{]}

Compute its dynamics:

\textbackslash{[}
\textbackslash dot\{z\}\_2 = \textbackslash ddot\{q\} - \textbackslash ddot\{q\}\_d + K\_1\textbackslash dot\{z\}\_1
= M\^{}\{-1\}(\textbackslash tau - C\textbackslash dot\{q\} - G) - \textbackslash ddot\{q\}\_d + K\_1(\textbackslash dot\{q\} - \textbackslash dot\{q\}\_d)
\textbackslash{]}

Augment the Lyapunov function (using the inertia matrix for natural energy weighting):

\textbackslash{[}
V\_2 = \textbackslash frac\{1\}\{2\}\textbackslash,z\_1\^{}T z\_1 + \textbackslash frac\{1\}\{2\}\textbackslash,z\_2\^{}T M(q)\textbackslash,z\_2
\textbackslash{]}

The derivative of \textbackslash(V\_2\textbackslash) is:

\textbackslash{[}
\textbackslash dot\{V\}\_2 = z\_1\^{}T\textbackslash dot\{z\}\_1 + z\_2\^{}T M\textbackslash dot\{z\}\_2 + \textbackslash frac\{1\}\{2\}\textbackslash,z\_2\^{}T\textbackslash dot\{M\}\textbackslash,z\_2
\textbackslash{]}

Substituting and using the skew-symmetry property \textbackslash(z\_2\^{}T(\textbackslash dot\{M\} - 2C)z\_2 = 0\textbackslash):

\textbackslash{[}
\textbackslash dot\{V\}\_2 = -z\_1\^{}T K\_1\textbackslash,z\_1 + z\_1\^{}T z\_2 + z\_2\^{}T\textbackslash bigl{[}\textbackslash tau - C\textbackslash dot\{q\} - G
- M(\textbackslash ddot\{q\}\_d - K\_1\textbackslash dot\{z\}\_1)\textbackslash bigr{]} + \textbackslash frac\{1\}\{2\}\textbackslash,z\_2\^{}T\textbackslash dot\{M\}\textbackslash,z\_2
\textbackslash{]}

Using the skew-symmetry property and simplifying:

\textbackslash{[}
\textbackslash dot\{V\}\_2 = -z\_1\^{}T K\_1\textbackslash,z\_1 + z\_1\^{}T z\_2 + z\_2\^{}T\textbackslash bigl{[}\textbackslash tau - C\textbackslash alpha\_1 - G
- M(\textbackslash ddot\{q\}\_d - K\_1\textbackslash dot\{z\}\_1)\textbackslash bigr{]}
\textbackslash{]}

To achieve \textbackslash(\textbackslash dot\{V\}\_2 = -z\_1\^{}T K\_1 z\_1 - z\_2\^{}T K\_2 z\_2\textbackslash), choose:

\textbf{Backstepping Torque Control Law:}\\
\strut \\
\textbackslash{[}
\textbackslash tau = M(q)\textbackslash bigl{[}\textbackslash ddot\{q\}\_d - K\_1\textbackslash dot\{z\}\_1 - K\_2 z\_2\textbackslash bigr{]}
+ C(q,\textbackslash dot\{q\})\textbackslash,\textbackslash alpha\_1 + G(q) - z\_1
\textbackslash{]}

where \textbackslash(K\_2 = k\_2 I \textbackslash succ 0\textbackslash). Expanding the terms:

\textbackslash{[}
\textbackslash tau = M(q)\textbackslash bigl{[}\textbackslash ddot\{q\}\_d - K\_1(\textbackslash dot\{q\} - \textbackslash dot\{q\}\_d) - K\_2(\textbackslash dot\{q\} - \textbackslash dot\{q\}\_d + K\_1(q-q\_d))\textbackslash bigr{]}
+ C(q,\textbackslash dot\{q\})(\textbackslash dot\{q\}\_d - K\_1(q-q\_d)) + G(q) - (q - q\_d)
\textbackslash{]}

\textbf{Stability result:}

\textbackslash{[}
\textbackslash dot\{V\}\_2 = -z\_1\^{}T K\_1\textbackslash,z\_1 - z\_2\^{}T K\_2\textbackslash,z\_2 \textbackslash le 0
\textbackslash{]}

Both \textbackslash(z\_1 \textbackslash to 0\textbackslash) and \textbackslash(z\_2 \textbackslash to 0\textbackslash) asymptotically, which means \textbackslash(q \textbackslash to q\_d\textbackslash) and \textbackslash(\textbackslash dot\{q\} \textbackslash to \textbackslash dot\{q\}\_d\textbackslash).

\subsubsection{4.2 MATLAB Implementation}\label{matlab-implementation}}

\% Backstepping controller for 2-DOF manipulator
function tau = backstepping\_manipulator(t, q, dq, params)
\% Desired trajectory
qd = {[}0.5*sin(t); 0.5*cos(t){]};
dqd = {[}0.5*cos(t); -0.5*sin(t){]};
ddqd = {[}-0.5*sin(t); -0.5*cos(t){]};
\% Errors
z1 = q - qd; \% Position error
dz1 = dq - dqd; \% Velocity of position error
alpha1 = dqd - params.k1 * z1; \% Virtual control (desired velocity)
z2 = dq - alpha1; \% Velocity error (= dz1 + k1*z1)
\% Dynamics matrices
{[}M, C, G{]} = manipulator\_dynamics\_matrices(q, dq, params);
\% Backstepping control law
tau = M * (ddqd - params.k1*dz1 - params.k2*z2) ...
+ C * alpha1 + G - z1;
end

\textbf{Implementation note:} The term \textbackslash(-z\_1\textbackslash) in the control law is the
"cross-coupling cancellation" term that arises from the backstepping procedure.
Without it, the position and velocity errors would remain coupled and stability
cannot be guaranteed by our Lyapunov function.

\subsection{Section 5 --- Part 3: Backstepping for Mobile Robot}\label{section-5-part-3-backstepping-for-mobile-robot}}

\subsubsection{5.1 Unicycle Model and Error Coordinates}\label{unicycle-model-and-error-coordinates}}

The unicycle (differential-drive) mobile robot has the kinematic model:

\textbackslash{[}
\textbackslash begin\{aligned\}
\textbackslash dot\{x\} \&= v\textbackslash cos\textbackslash theta \textbackslash\textbackslash{}
\textbackslash dot\{y\} \&= v\textbackslash sin\textbackslash theta \textbackslash\textbackslash{}
\textbackslash dot\{\textbackslash theta\} \&= \textbackslash omega
\textbackslash end\{aligned\}
\textbackslash{]}

where \textbackslash(v\textbackslash) is the linear velocity and \textbackslash(\textbackslash omega\textbackslash) is the angular velocity (our two control inputs).

Given a reference trajectory \textbackslash((x\_r(t), y\_r(t), \textbackslash theta\_r(t))\textbackslash) generated by
\textbackslash(\textbackslash dot\{x\}\_r = v\_r\textbackslash cos\textbackslash theta\_r\textbackslash), \textbackslash(\textbackslash dot\{y\}\_r = v\_r\textbackslash sin\textbackslash theta\_r\textbackslash), \textbackslash(\textbackslash dot\{\textbackslash theta\}\_r = \textbackslash omega\_r\textbackslash),
define the \textbf{body-frame error}:

\textbackslash{[}
\textbackslash begin\{bmatrix\} e\_1 \textbackslash\textbackslash{} e\_2 \textbackslash\textbackslash{} e\_3 \textbackslash end\{bmatrix\}
= \textbackslash begin\{bmatrix\} \textbackslash cos\textbackslash theta \& \textbackslash sin\textbackslash theta \& 0 \textbackslash\textbackslash{} -\textbackslash sin\textbackslash theta \& \textbackslash cos\textbackslash theta \& 0 \textbackslash\textbackslash{} 0 \& 0 \& 1 \textbackslash end\{bmatrix\}
\textbackslash begin\{bmatrix\} x\_r - x \textbackslash\textbackslash{} y\_r - y \textbackslash\textbackslash{} \textbackslash theta\_r - \textbackslash theta \textbackslash end\{bmatrix\}
\textbackslash{]}

The error dynamics are:

\textbackslash{[}
\textbackslash begin\{aligned\}
\textbackslash dot\{e\}\_1 \&= \textbackslash omega\textbackslash, e\_2 - v + v\_r\textbackslash cos e\_3 \textbackslash\textbackslash{}
\textbackslash dot\{e\}\_2 \&= -\textbackslash omega\textbackslash, e\_1 + v\_r\textbackslash sin e\_3 \textbackslash\textbackslash{}
\textbackslash dot\{e\}\_3 \&= \textbackslash omega\_r - \textbackslash omega
\textbackslash end\{aligned\}
\textbackslash{]}

\textbf{Figure 6:} Differential-drive mobile robot (top-down view) with body-frame errors and the backstepping cascade: position error drives heading control, which drives velocity control.

\subsubsection{5.2 Backstepping Design}\label{backstepping-design}}

\paragraph{Step 1 --- Stabilise the Lateral/Heading Subsystem}\label{step-1-stabilise-the-lateralheading-subsystem}}

Consider the \textbackslash((e\_2, e\_3)\textbackslash) subsystem first. Note that \textbackslash(e\_2\textbackslash) is driven by \textbackslash(v\_r \textbackslash sin e\_3\textbackslash).
Think of \textbackslash(e\_3\textbackslash) as a virtual input to steer \textbackslash(e\_2\textbackslash). If \textbackslash(e\_3\textbackslash) could be freely chosen, pick:

\textbackslash{[}
e\_3\^{}* = \textbackslash arctan(-k\_3\textbackslash, e\_2)
\textbackslash{]}

(This makes \textbackslash(v\_r \textbackslash sin e\_3\^{}* \textbackslash approx -k\_3 v\_r e\_2\textbackslash) for small \textbackslash(e\_2\textbackslash), driving \textbackslash(e\_2 \textbackslash to 0\textbackslash).)

However, \textbackslash(e\_3\textbackslash) is a state, not an input. Define the backstepping error:

\textbackslash{[}
\textbackslash bar\{e\}\_3 = e\_3 - e\_3\^{}*
\textbackslash{]}

\paragraph{Step 2 --- Design the Control Inputs \textbackslash((v, \textbackslash omega)\textbackslash)}\label{step-2-design-the-control-inputs-v-omega}}

Choose the Lyapunov function:

\textbackslash{[}
V = \textbackslash frac\{1\}\{2\}e\_1\^{}2 + \textbackslash frac\{1\}\{2\}e\_2\^{}2 + \textbackslash frac\{1\}\{2\}\textbackslash bar\{e\}\_3\^{}2
\textbackslash{]}

For the longitudinal error \textbackslash(e\_1\textbackslash), choose the linear velocity:

\textbackslash{[}
v = v\_r\textbackslash cos e\_3 + k\_1\textbackslash,e\_1
\textbackslash{]}

This gives \textbackslash(\textbackslash dot\{e\}\_1 = -k\_1 e\_1 + \textbackslash omega e\_2\textbackslash).

For the angular velocity, we need \textbackslash(\textbackslash omega\textbackslash) to simultaneously:
(a) cancel the coupling \textbackslash(\textbackslash omega e\_2\textbackslash) term in \textbackslash(\textbackslash dot\{e\}\_1\textbackslash),
(b) drive \textbackslash(\textbackslash bar\{e\}\_3 \textbackslash to 0\textbackslash), and
(c) handle the \textbackslash(e\_2\textbackslash) dynamics.

The full angular velocity controller is:

\textbackslash{[}
\textbackslash omega = \textbackslash omega\_r + k\_2\textbackslash,v\_r\textbackslash,e\_2 + k\_3\textbackslash,\textbackslash bar\{e\}\_3
\textbackslash{]}

where \textbackslash(k\_2\textbackslash) handles the lateral error coupling with \textbackslash(v\_r \textbackslash sin e\_3\textbackslash), and \textbackslash(k\_3\textbackslash) drives the heading backstepping error to zero.

\textbf{Backstepping Mobile Robot Controller:}\\
\strut \\
\textbackslash{[}
v = v\_r \textbackslash cos e\_3 + k\_1\textbackslash,e\_1
\textbackslash{]}
\textbackslash{[}
\textbackslash omega = \textbackslash omega\_r + k\_2\textbackslash,v\_r\textbackslash,e\_2 + k\_3\textbackslash,\textbackslash sin e\_3
\textbackslash{]}

\textbf{Stability analysis.} With this controller and \textbackslash(V = \textbackslash frac\{1\}\{2\}(e\_1\^{}2 + e\_2\^{}2 + e\_3\^{}2)\textbackslash):

\textbackslash{[}
\textbackslash dot\{V\} = e\_1\textbackslash dot\{e\}\_1 + e\_2\textbackslash dot\{e\}\_2 + e\_3\textbackslash dot\{e\}\_3
\textbackslash{]}
\textbackslash{[}
= e\_1(-k\_1 e\_1 + \textbackslash omega e\_2) + e\_2(-\textbackslash omega e\_1 + v\_r \textbackslash sin e\_3)
+ e\_3(\textbackslash omega\_r - \textbackslash omega)
\textbackslash{]}

The cross-terms \textbackslash(e\_1 \textbackslash omega e\_2\textbackslash) and \textbackslash(-e\_2 \textbackslash omega e\_1\textbackslash) cancel. Substituting the controller:

\textbackslash{[}
\textbackslash dot\{V\} = -k\_1 e\_1\^{}2 + e\_2 v\_r \textbackslash sin e\_3 + e\_3(\textbackslash omega\_r - \textbackslash omega\_r - k\_2 v\_r e\_2 - k\_3 \textbackslash sin e\_3)
\textbackslash{]}
\textbackslash{[}
= -k\_1 e\_1\^{}2 + v\_r e\_2 \textbackslash sin e\_3 - k\_2 v\_r e\_2 e\_3 - k\_3 e\_3 \textbackslash sin e\_3
\textbackslash{]}

For \textbackslash(v\_r \textgreater{} 0\textbackslash) and using the fact that \textbackslash(\textbackslash frac\{\textbackslash sin e\_3\}\{e\_3\}\textbackslash) is bounded, with appropriate choice of \textbackslash(k\_2\textbackslash), we get \textbackslash(\textbackslash dot\{V\} \textbackslash le -k\_1 e\_1\^{}2 - k\_3 e\_3 \textbackslash sin e\_3 \textbackslash le 0\textbackslash). Convergence of all errors follows by Barbalat\textquotesingle s lemma (since \textbackslash(v\_r(t) \textgreater{} 0\textbackslash)).

\subsubsection{5.3 MATLAB Implementation}\label{matlab-implementation-1}}

\% Backstepping controller for unicycle mobile robot
function {[}v, omega{]} = backstepping\_mobile\_robot(e1, e2, e3, v\_r, omega\_r, params)
\% e1: longitudinal error (body frame)
\% e2: lateral error (body frame)
\% e3: heading error
v = v\_r * cos(e3) + params.k1 * e1;
omega = omega\_r + params.k2 * v\_r * e2 + params.k3 * sin(e3);
end

\subsection{Section 6 --- Part 4: Backstepping for Quadrotor Drone}\label{section-6-part-4-backstepping-for-quadrotor-drone}}

\subsubsection{6.1 Quadrotor Cascade Structure}\label{quadrotor-cascade-structure}}

The quadrotor has a natural four-level cascade:

\textbackslash{[}
\textbackslash underbrace\{\textbackslash text\{Position\}\}\_\{\textbackslash text\{Level 1\}\}
\textbackslash;\textbackslash xrightarrow\{\textbackslash text\{desired accel.\}\}\textbackslash;
\textbackslash underbrace\{\textbackslash text\{Velocity\}\}\_\{\textbackslash text\{Level 2\}\}
\textbackslash;\textbackslash xrightarrow\{\textbackslash text\{desired angles\}\}\textbackslash;
\textbackslash underbrace\{\textbackslash text\{Orientation\}\}\_\{\textbackslash text\{Level 3\}\}
\textbackslash;\textbackslash xrightarrow\{\textbackslash text\{desired rates\}\}\textbackslash;
\textbackslash underbrace\{\textbackslash text\{Angular Rate\}\}\_\{\textbackslash text\{Level 4\}\}
\textbackslash;\textbackslash xrightarrow\{\textbackslash quad\}\textbackslash;
(T, \textbackslash boldsymbol\{\textbackslash tau\})
\textbackslash{]}

The simplified quadrotor dynamics (with rotation matrix \textbackslash(R\textbackslash)):

\textbackslash{[}
\textbackslash begin\{aligned\}
\textbackslash dot\{\textbackslash mathbf\{p\}\} \&= \textbackslash mathbf\{v\} \textbackslash\textbackslash{[}4pt{]}
m\textbackslash dot\{\textbackslash mathbf\{v\}\} \&= -mg\textbackslash,\textbackslash mathbf\{e\}\_3 + T\textbackslash,R\textbackslash,\textbackslash mathbf\{e\}\_3 \textbackslash\textbackslash{[}4pt{]}
\textbackslash dot\{R\} \&= R\textbackslash,{[}\textbackslash boldsymbol\{\textbackslash omega\}{]}\_\textbackslash times \textbackslash\textbackslash{[}4pt{]}
J\textbackslash dot\{\textbackslash boldsymbol\{\textbackslash omega\}\} \&= -\textbackslash boldsymbol\{\textbackslash omega\} \textbackslash times J\textbackslash boldsymbol\{\textbackslash omega\} + \textbackslash boldsymbol\{\textbackslash tau\}\_b
\textbackslash end\{aligned\}
\textbackslash{]}

where \textbackslash(\textbackslash mathbf\{p\} = {[}x,y,z{]}\^{}T\textbackslash) is position, \textbackslash(\textbackslash mathbf\{v\}\textbackslash) is velocity, \textbackslash(T\textbackslash) is total thrust,
\textbackslash(R \textbackslash in SO(3)\textbackslash) is the rotation matrix, \textbackslash(\textbackslash boldsymbol\{\textbackslash omega\}\textbackslash) is the body angular velocity,
\textbackslash(J\textbackslash) is the inertia tensor, and \textbackslash(\textbackslash boldsymbol\{\textbackslash tau\}\_b\textbackslash) are the body torques.

\textbf{Figure 7:} Quadrotor cascade backstepping architecture. The position controller generates a desired attitude (virtual control); the attitude controller generates the actual motor commands.

\subsubsection{6.2 Backstepping Design --- Outer Loop (Position Control)}\label{backstepping-design-outer-loop-position-control}}

\paragraph{Step 1 --- Position Error}\label{step-1-position-error}}

Define position error \textbackslash(z\_1 = \textbackslash mathbf\{p\} - \textbackslash mathbf\{p\}\_d\textbackslash). Then \textbackslash(\textbackslash dot\{z\}\_1 = \textbackslash mathbf\{v\} - \textbackslash dot\{\textbackslash mathbf\{p\}\}\_d\textbackslash).

Virtual control (desired velocity):

\textbackslash{[}
\textbackslash boldsymbol\{\textbackslash alpha\}\_1 = \textbackslash dot\{\textbackslash mathbf\{p\}\}\_d - K\_\{p1\}\textbackslash,z\_1, \textbackslash qquad K\_\{p1\} = k\_\{p1\}\textbackslash,I\_3 \textbackslash succ 0
\textbackslash{]}

\paragraph{Step 2 --- Velocity Error}\label{step-2-velocity-error}}

Define velocity error \textbackslash(z\_2 = \textbackslash mathbf\{v\} - \textbackslash boldsymbol\{\textbackslash alpha\}\_1\textbackslash). Its dynamics:

\textbackslash{[}
\textbackslash dot\{z\}\_2 = \textbackslash dot\{\textbackslash mathbf\{v\}\} - \textbackslash dot\{\textbackslash boldsymbol\{\textbackslash alpha\}\}\_1
= -g\textbackslash mathbf\{e\}\_3 + \textbackslash frac\{T\}\{m\}R\textbackslash mathbf\{e\}\_3 - \textbackslash ddot\{\textbackslash mathbf\{p\}\}\_d + K\_\{p1\}\textbackslash dot\{z\}\_1
\textbackslash{]}

Choose the Lyapunov function \textbackslash(V = \textbackslash frac\{1\}\{2\}z\_1\^{}T z\_1 + \textbackslash frac\{1\}\{2\}z\_2\^{}T z\_2\textbackslash).
To make \textbackslash(\textbackslash dot\{V\} = -k\_\{p1\}\textbackslash\textbar z\_1\textbackslash\textbar\^{}2 - k\_\{p2\}\textbackslash\textbar z\_2\textbackslash\textbar\^{}2\textbackslash), we need the
\textbf{desired acceleration vector}:

\textbackslash{[}
\textbackslash mathbf\{a\}\_\{\textbackslash text\{des\}\} = \textbackslash ddot\{\textbackslash mathbf\{p\}\}\_d - K\_\{p1\}\textbackslash dot\{z\}\_1 - z\_1 - K\_\{p2\}\textbackslash,z\_2 + g\textbackslash mathbf\{e\}\_3
\textbackslash{]}

The thrust magnitude and desired orientation are extracted from \textbackslash(\textbackslash mathbf\{a\}\_\{\textbackslash text\{des\}\}\textbackslash):

\textbackslash{[}
T = m\textbackslash,\textbackslash\textbar\textbackslash mathbf\{a\}\_\{\textbackslash text\{des\}\}\textbackslash\textbar{}
\textbackslash{]}
\textbackslash{[}
\textbackslash mathbf\{z\}\_\{B,\textbackslash text\{des\}\} = \textbackslash frac\{\textbackslash mathbf\{a\}\_\{\textbackslash text\{des\}\}\}\{\textbackslash\textbar\textbackslash mathbf\{a\}\_\{\textbackslash text\{des\}\}\textbackslash\textbar\}
\textbackslash quad \textbackslash Longrightarrow \textbackslash quad
R\_\{\textbackslash text\{des\}\} \textbackslash text\{ constructed from \} \textbackslash mathbf\{z\}\_\{B,\textbackslash text\{des\}\} \textbackslash text\{ and desired yaw\}
\textbackslash{]}

\subsubsection{6.3 Backstepping Design --- Inner Loop (Attitude Control)}\label{backstepping-design-inner-loop-attitude-control}}

\paragraph{Step 3 --- Attitude Error}\label{step-3-attitude-error}}

Using Euler angles \textbackslash((\textbackslash phi, \textbackslash theta, \textbackslash psi)\textbackslash) for a simplified treatment, define:

\textbackslash{[}
z\_3 = \textbackslash boldsymbol\{\textbackslash Phi\} - \textbackslash boldsymbol\{\textbackslash Phi\}\_\{\textbackslash text\{des\}\}
= \textbackslash begin\{bmatrix\} \textbackslash phi - \textbackslash phi\_d \textbackslash\textbackslash{} \textbackslash theta - \textbackslash theta\_d \textbackslash\textbackslash{} \textbackslash psi - \textbackslash psi\_d \textbackslash end\{bmatrix\}
\textbackslash{]}

Virtual control (desired angular rate):

\textbackslash{[}
\textbackslash boldsymbol\{\textbackslash alpha\}\_3 = \textbackslash dot\{\textbackslash boldsymbol\{\textbackslash Phi\}\}\_\{\textbackslash text\{des\}\} - K\_\{a1\}\textbackslash,z\_3
\textbackslash{]}

\paragraph{Step 4 --- Angular Rate Error and Torque Design}\label{step-4-angular-rate-error-and-torque-design}}

Define \textbackslash(z\_4 = \textbackslash boldsymbol\{\textbackslash omega\} - \textbackslash boldsymbol\{\textbackslash alpha\}\_3\textbackslash). Using the Euler equation
\textbackslash(J\textbackslash dot\{\textbackslash boldsymbol\{\textbackslash omega\}\} = -\textbackslash boldsymbol\{\textbackslash omega\} \textbackslash times J\textbackslash boldsymbol\{\textbackslash omega\} + \textbackslash boldsymbol\{\textbackslash tau\}\_b\textbackslash):

\textbf{Backstepping Attitude Torque:}\\
\strut \\
\textbackslash{[}
\textbackslash boldsymbol\{\textbackslash tau\}\_b = J\textbackslash bigl{[}\textbackslash dot\{\textbackslash boldsymbol\{\textbackslash alpha\}\}\_3 - K\_\{a2\}\textbackslash,z\_4\textbackslash bigr{]}
+ \textbackslash boldsymbol\{\textbackslash omega\} \textbackslash times J\textbackslash boldsymbol\{\textbackslash omega\} - z\_3
\textbackslash{]}

This yields the overall Lyapunov function:

\textbackslash{[}
V = \textbackslash frac\{1\}\{2\}\textbackslash bigl(\textbackslash\textbar z\_1\textbackslash\textbar\^{}2 + \textbackslash\textbar z\_2\textbackslash\textbar\^{}2 + \textbackslash\textbar z\_3\textbackslash\textbar\^{}2 + z\_4\^{}T J\textbackslash,z\_4\textbackslash bigr)
\textbackslash{]}
\textbackslash{[}
\textbackslash dot\{V\} = -k\_\{p1\}\textbackslash\textbar z\_1\textbackslash\textbar\^{}2 - k\_\{p2\}\textbackslash\textbar z\_2\textbackslash\textbar\^{}2 - k\_\{a1\}\textbackslash\textbar z\_3\textbackslash\textbar\^{}2 - k\_\{a2\}\textbackslash,z\_4\^{}T K\_\{a2\}\textbackslash,z\_4 \textbackslash le 0
\textbackslash{]}

\subsubsection{6.4 MATLAB Implementation}\label{matlab-implementation-2}}

\% Backstepping controller for quadrotor
function {[}T, tau\_b{]} = backstepping\_quadrotor(state, ref, params)
\% State: {[}x y z vx vy vz phi theta psi p q r{]}
pos = state(1:3); vel = state(4:6);
phi = state(7); theta = state(8); psi = state(9);
omega\_b = state(10:12);
\% Reference
pos\_d = ref.pos; vel\_d = ref.vel; acc\_d = ref.acc; psi\_d = ref.psi;
\% -\/-\/- Outer loop: Position backstepping -\/-\/-
z1 = pos - pos\_d;
dz1 = vel - vel\_d;
alpha1 = vel\_d - params.kp1 * z1; \% Virtual velocity
z2 = vel - alpha1; \% Velocity error
\% Desired acceleration
a\_des = acc\_d - params.kp1*dz1 - z1 - params.kp2*z2 + {[}0;0;params.g{]};
\% Thrust
T = params.m * norm(a\_des);
\% Desired orientation from a\_des
z\_B\_des = a\_des / norm(a\_des);
phi\_d = asin(z\_B\_des(1)*sin(psi\_d) - z\_B\_des(2)*cos(psi\_d));
theta\_d = atan2(z\_B\_des(1)*cos(psi\_d)+z\_B\_des(2)*sin(psi\_d), z\_B\_des(3));
\% -\/-\/- Inner loop: Attitude backstepping -\/-\/-
Phi = {[}phi; theta; psi{]};
Phi\_d = {[}phi\_d; theta\_d; psi\_d{]};
z3 = Phi - Phi\_d;
alpha3 = -params.ka1 * z3; \% Desired angular rate
z4 = omega\_b - alpha3;
\% Body torques
J = diag({[}params.Ixx, params.Iyy, params.Izz{]});
tau\_b = J * (-params.ka1*(-params.ka1*z3 - z4*0) - params.ka2*z4) ...
+ cross(omega\_b, J*omega\_b) - z3;
end

\textbf{Practical note:} In real implementations, the desired angles \textbackslash((\textbackslash phi\_d, \textbackslash theta\_d)\textbackslash)
must be saturated (typically to \textbackslash(\textbackslash pm 30\^{}\textbackslash circ\textbackslash)) to prevent aggressive manoeuvres.
The thrust \textbackslash(T\textbackslash) must also be clamped to \textbackslash({[}0, T\_\{\textbackslash max\}{]}\textbackslash).

\subsection{Section 7 --- Hands-on Activities \& Quiz}\label{section-7-hands-on-activities-quiz}}

\subsubsection{Activity 1: Gain Tuning for Manipulator Backstepping (10 min)}\label{activity-1-gain-tuning-for-manipulator-backstepping-10-min}}

Run \texttt{backstepping\_control\_simulation.m} and observe the manipulator plots.
Then modify the gains and re-run:

\begin{enumerate}
\tightlist
\item
  Set \textbackslash(k\_1 = 50, \textbackslash; k\_2 = 5\textbackslash) (position gain much larger than velocity gain).
\item
  Set \textbackslash(k\_1 = 5, \textbackslash; k\_2 = 50\textbackslash) (velocity gain much larger than position gain).
\item
  Set \textbackslash(k\_1 = k\_2 = 25\textbackslash) (balanced gains).
\end{enumerate}

\textbf{Expected observations:}

\begin{itemize}
\tightlist
\item
  \textbf{\textbackslash(k\_1 \textbackslash gg k\_2\textbackslash):} Aggressive position correction generates large desired velocity changes.
  The velocity loop cannot track them fast enough, causing oscillatory or even unstable behaviour.
  The virtual control \textbackslash(\textbackslash alpha\_1\textbackslash) demands more than the torque loop can deliver.
\item
  \textbf{\textbackslash(k\_2 \textbackslash gg k\_1\textbackslash):} The velocity error \textbackslash(z\_2\textbackslash) converges quickly, but the position
  error \textbackslash(z\_1\textbackslash) converges slowly because \textbackslash(k\_1\textbackslash) is small. The system is overdamped
  with sluggish position tracking.
\item
  \textbf{\textbackslash(k\_1 \textbackslash approx k\_2\textbackslash):} Balanced convergence rates for both loops. This is
  generally the best choice for smooth, fast tracking.
\end{itemize}

\subsubsection{Activity 2: Comparative Study (15 min)}\label{activity-2-comparative-study-15-min}}

The MATLAB file includes a comparison of backstepping vs. computed torque vs. robust control
on the manipulator with model uncertainty (20\% mass error). Observe:

\begin{itemize}
\tightlist
\item
  Which controller has the smallest steady-state error?
\item
  Which controller has the smoothest control signal?
\item
  Which controller uses the least energy (smallest \textbackslash(\textbackslash int \textbackslash\textbar\textbackslash tau\textbackslash\textbar\^{}2 \textbackslash, dt\textbackslash))?
\end{itemize}

\paragraph{Quiz --- Virtual Control Design}\label{quiz-virtual-control-design}}

\textbf{Q1:} Given the system \textbackslash(\textbackslash dot\{x\}\_1 = x\_1\^{}2 + x\_2\textbackslash), \textbackslash(\textbackslash dot\{x\}\_2 = x\_2 + u\textbackslash),
design a backstepping controller to stabilise \textbackslash(x\_1 = 0\textbackslash).

\emph{Answer sketch:}

\begin{itemize}
\tightlist
\item
  Step 1: \textbackslash(z\_1 = x\_1\textbackslash), \textbackslash(\textbackslash dot\{z\}\_1 = x\_1\^{}2 + x\_2\textbackslash).
  Choose \textbackslash(\textbackslash alpha\_1 = -x\_1\^{}2 - k\_1 x\_1\textbackslash) so that \textbackslash(\textbackslash dot\{V\}\_1 = -k\_1 x\_1\^{}2 + x\_1 z\_2\textbackslash).
\item
  Step 2: \textbackslash(z\_2 = x\_2 - \textbackslash alpha\_1\textbackslash), \textbackslash(\textbackslash dot\{z\}\_2 = x\_2 + u - (2x\_1)(x\_1\^{}2 + x\_2) + k\_1(x\_1\^{}2 + x\_2)\textbackslash).
  Choose \textbackslash(u = -x\_2 + \textbackslash dot\{\textbackslash alpha\}\_1 - x\_1 - k\_2 z\_2\textbackslash).
\end{itemize}

\textbf{Q2:} In the two-step backstepping controller, what happens if we omit the
cross-coupling term \textbackslash(-z\_1\textbackslash) from the torque law? Is stability still guaranteed?

\emph{Answer:} Without \textbackslash(-z\_1\textbackslash), the Lyapunov derivative becomes
\textbackslash(\textbackslash dot\{V\}\_2 = -k\_1 z\_1\^{}2 - k\_2 z\_2\^{}2 + z\_1\^{}T z\_2\textbackslash). The cross-term \textbackslash(z\_1\^{}T z\_2\textbackslash) can be
positive, so \textbackslash(\textbackslash dot\{V\}\_2\textbackslash) is not guaranteed negative semi-definite. Stability is only ensured
if \textbackslash(k\_1 k\_2 \textgreater{} \textbackslash frac\{1\}\{4\}\textbackslash) (by Young\textquotesingle s inequality), but the stability margin is reduced.
The systematic backstepping design avoids this issue entirely.

\textbf{Q3:} Why does the quadrotor need four backstepping steps while the
manipulator only needs two?

\emph{Answer:} The quadrotor has a four-level cascade:
position → velocity → orientation → angular rate, each acting as a virtual control
for the previous level. The manipulator only has two levels (position → velocity) because
the torque directly enters the velocity dynamics via \textbackslash(M\^{}\{-1\}\textbackslash tau\textbackslash). The number of backstepping
steps equals the relative degree from the output to the input.

\subsection{Section 8 --- Summary: Comparison of All Five Control Methods}\label{section-8-summary-comparison-of-all-five-control-methods}}

\begin{longtable}[]{@{}llllll@{}}
\toprule\noalign{}
Aspect & Computed Torque (Wk1) & MPC (Wk2) & Adaptive (Wk3) & Robust (Wk4) & Backstepping (Wk5) \\
\midrule\noalign{}
\endhead
\bottomrule\noalign{}
\endlastfoot
\textbf{Core idea} & Cancel nonlinearity, impose linear error dynamics & Optimise predicted trajectory over finite horizon & On-line parameter estimation + control & Dominate worst-case bounded uncertainty & Recursive Lyapunov design through cascade structure \\
\textbf{Model requirement} & Exact model needed & Prediction model needed & Structure known, parameters unknown & Nominal model + uncertainty bound & Strict-feedback structure required \\
\textbf{Lyapunov function} & Implicit (linear closed-loop) & Cost function (not standard Lyapunov) & \textbackslash(V = \textbackslash frac\{1\}\{2\}s\^{}T M s + \textbackslash frac\{1\}\{2\}\textbackslash tilde\{\textbackslash theta\}\^{}T\textbackslash Gamma\^{}\{-1\}\textbackslash tilde\{\textbackslash theta\}\textbackslash) & \textbackslash(V = \textbackslash frac\{1\}\{2\}s\^{}T M s\textbackslash) & \textbackslash(V = \textbackslash sum \textbackslash frac\{1\}\{2\}z\_i\^{}T W\_i z\_i\textbackslash) (constructed step-by-step) \\
\textbf{Handles constraints} & No & Yes (core feature) & No (standard form) & No & No (standard form) \\
\textbf{Robustness} & Poor (sensitive to model errors) & Moderate (depends on model accuracy) & Good (adapts to parametric uncertainty) & Excellent (designed for uncertainty) & Good (Lyapunov-guaranteed; can add adaptation) \\
\textbf{Computation} & Low (closed-form) & High (on-line optimisation) & Moderate (adaptation law update) & Low (closed-form) & Moderate (derivatives of virtual controls) \\
\textbf{Key advantage} & Simple, well-understood & Handles constraints and preview & Self-tuning for unknown parameters & Guaranteed performance under uncertainty & Systematic design with constructive stability proof \\
\textbf{Key limitation} & No robustness to model error & Computational cost, model dependence & Needs persistent excitation, structured uncertainty & Conservative (worst-case design), chattering & Complexity explosion in \textbackslash(\textbackslash dot\{\textbackslash alpha\}\_i\textbackslash) for high-order systems \\
\textbf{Control law form} & \textbackslash(\textbackslash tau = M(q){[}\textbackslash ddot\{q\}\_d + K\_D\textbackslash dot\{e\} + K\_P e{]} + C\textbackslash dot\{q\} + G\textbackslash) & \textbackslash(u\^{}* = \textbackslash arg\textbackslash min\_\{u\} \textbackslash sum\_\{k\} \textbackslash\textbar x\_k - x\_r\textbackslash\textbar\^{}2\_Q + \textbackslash\textbar u\_k\textbackslash\textbar\^{}2\_R\textbackslash) & \textbackslash(\textbackslash tau = Y(q,\textbackslash dot\{q\},\textbackslash dot\{q\}\_r,\textbackslash ddot\{q\}\_r)\textbackslash hat\{\textbackslash theta\} - K\_D s\textbackslash) & \textbackslash(\textbackslash tau = \textbackslash hat\{M\}\textbackslash ddot\{q\}\_r + \textbackslash hat\{C\}\textbackslash dot\{q\}\_r + \textbackslash hat\{G\} - K\_D s + \textbackslash rho\textbackslash frac\{s\}\{\textbackslash\textbar s\textbackslash\textbar\}\textbackslash) & \textbackslash(\textbackslash tau = M{[}\textbackslash ddot\{q\}\_d - K\_1\textbackslash dot\{z\}\_1 - K\_2 z\_2{]} + C\textbackslash alpha\_1 + G - z\_1\textbackslash) \\
\end{longtable}

\textbf{Figure 8:} Side-by-side comparison of feedback linearization (cancels all nonlinearities) versus backstepping (exploits useful nonlinearities with constructive Lyapunov guarantees).

\textbf{Key takeaway:} Backstepping provides a \emph{systematic, constructive} methodology
for designing stabilising controllers for nonlinear systems in strict-feedback form.
Unlike computed torque, it does not blindly cancel all nonlinearities.
Unlike robust control, it does not require conservative worst-case bounds.
Its main challenge --- the "explosion of terms" as \textbackslash(\textbackslash dot\{\textbackslash alpha\}\_i\textbackslash) grows with each step
--- can be addressed using \textbf{dynamic surface control} (a first-order filter on
each virtual control), which we may explore in future weeks.

\textbf{M.Sc. Control for Robots} --- University of West London --- AY 2025--26

Week 5: Backstepping Control --- Lecture Notes

These notes are for educational purposes.
Students are encouraged to verify all derivations independently.
